% Challenge dossier: VI.06.C1
\subsection*{Challenge VI.06.C1 --- The radical as an intersection of prime ideals}
\label{chal:vi-06-c1}

\paragraph{Problem.}
Let \(A\) be a commutative ring and \(I\subseteq A\) an ideal.  Prove the theorem
\[
\sqrt I
=
\bigcap_{\substack{\mathfrak p\supseteq I\\ \mathfrak p\text{ prime}}}\mathfrak p.
\]
You may use Zorn's lemma.

\ifdefined\FullProblemDossiers
\paragraph{Why this problem matters.}
This theorem says that prime ideals are sufficient to recover every radical closed
condition.  It is one of the deepest algebraic reasons that spectra use all prime ideals.

\paragraph{Useful ideas before starting.}
One inclusion is immediate because primes are radical.  For the reverse inclusion, if
\(f\notin\sqrt I\), localize or use a maximal ideal among ideals containing \(I\) and
disjoint from \(\{1,f,f^2,\dots\}\).

\paragraph{Guided hints.}
Let \(S=\{1,f,f^2,\dots\}\).  If \(f\notin\sqrt I\), then \(I\cap S=\varnothing\).  Use
Zorn's lemma to choose an ideal maximal with respect to containing \(I\) and avoiding \(S\),
and prove that it is prime.

\paragraph{Solution.}
Every prime \(\mathfrak p\supseteq I\) contains \(\sqrt I\), because if
\(a^n\in I\subseteq\mathfrak p\), primality gives \(a\in\mathfrak p\).  Hence
\[
\sqrt I\subseteq\bigcap_{\mathfrak p\supseteq I}\mathfrak p.
\]

For the reverse inclusion, take \(f\notin\sqrt I\).  Then no power of \(f\) lies in \(I\),
so \(I\cap S=\varnothing\) for \(S=\{1,f,f^2,\dots\}\).  By Zorn's lemma choose an ideal
\(\mathfrak p\) maximal among ideals containing \(I\) and disjoint from \(S\).  We claim it
is prime.  If \(ab\in\mathfrak p\) but neither \(a\) nor \(b\) lies in \(\mathfrak p\),
maximality implies both \(\mathfrak p+(a)\) and \(\mathfrak p+(b)\) meet \(S\).  Thus
\[
p_1+r a=f^m,\qquad p_2+s b=f^n
\]
for suitable \(p_1,p_2\in\mathfrak p\).  Multiplying shows \(f^{m+n}\in\mathfrak p\),
contradicting disjointness from \(S\).  Hence \(\mathfrak p\) is prime.  It contains \(I\)
but not \(f\).  Therefore \(f\) is not in the intersection of all primes over \(I\).

\paragraph{What happened mathematically.}
Prime ideals form enough geometric points to detect radical membership.  Equations vanish
on all prime points above \(I\) precisely when some power already lies in \(I\).

\paragraph{Extensions.}
Apply the theorem to \(I=(0)\) to recover the nilradical as the intersection of all prime
ideals.

\paragraph{Provenance.}
New canonical challenge based on the standard prime-existence theorem of commutative
algebra.
\fi
